\section{总结与延伸}

\subsection{全课总结表}
\begin{table}[H]
\centering
\begin{tabular}{p{2.9cm}p{4.3cm}p{4.5cm}}
\toprule
\textbf{主题} & \textbf{课程核心判断} & \textbf{实践动作} \\
\midrule
训练流程 & SFT $\rightarrow$ Reasoning RL $\rightarrow$ Agent RL & 按阶段构建数据和奖励，避免一步到位 \\
评估体系 & evaluation is the key & 闭集提速，开集保真，双轨并行 \\
成本约束 & 训练成本与推理成本必须分开管理 & 建立端到端成本仪表盘 \\
工具训练 & synthetic tools 可扩展能力边界 & 合成数据与真实回放联合训练 \\
系统瓶颈 & 长时任务常卡在内存与通信 & 异步调度、缓存与分层存储优化 \\
安全上线 & 风险控制贯穿训练到运行时 & 分阶段放权与失败恢复评估 \\
\bottomrule
\end{tabular}
\caption{Lecture 01 的可执行结论}
\end{table}

\subsection{关键结论回顾}
\begin{importantbox}{六条必须记住的结论}
\begin{enumerate}
    \item Agent 成功率取决于系统协同，而非单点模型分数。
    \item Reasoning RL 是连接文本能力与执行能力的关键桥梁。
    \item 数据与评估必须共设计，否则会互相污染。
    \item 成本分析必须覆盖训练、评估和推理全链路。
    \item 工具使用训练要关注决策质量，而不只是调用格式。
    \item 开集任务里，失败恢复能力比一次成功更重要。
\end{enumerate}
\end{importantbox}

\subsection{进一步阅读}
\begin{itemize}
    \item Berkeley RDI Agentic AI MOOC（Fall 2025）课程主页与讲义。
    \item Yann Dubois 关于 LLM 训练与评估的公开分享与论文。
    \item ReAct、SWE-bench、WebArena、OSWorld 等 Agent 任务与评测工作。
    \item 工具使用与长任务强化学习相关实践（tool simulation、async rollout）。
    \item 推理模型与 Agent 模型结合的后训练研究方向。
\end{itemize}

\end{document}
